\scanpage{13}
\mspage{4}
Je considère maintenant la sous catégorie ${\unsure{\mathscr B|\mathcal B|B}}_0$ de $\widehat M\supset M$, formée des objets $X$ de la forme suivante :
\[
\left(\begin{aligned}
&X\text{ est }\deleted{\unsure{\varinjlim|\underset{\longrightarrow}{\text{limite}}}}\inserted{\text{sup}}\text{ de ses sous-objets isomorphes à des}\\
&\deleted{\text{\unsure{les|ses|des}}}\qquad\text{objets de }M\text{ (modèles)}
\end{aligned}\right.
\]
% slanted
\margin{Cela \unsure{implique|signifie} que \unsure{tt|$\forall$} hom\unsure{,} \deleted{de} $S\unsure{-|\to}X$, où $S\in\operatorname{Ob}M$, est un mono}

Si $K,L\subset X$ sont des sous-objets modèles, alors $K\cap L$ est sup \unsure{des|de} sous-objets de $K\cap L$ qui sont des modèles

Soit Donc $\widetilde X$ l'ens des \deleted{\unsure{obj}} sous-objets de $X$ qui sont des modèles. Alors $\widetilde X$. est un ens ordonné \inserted{$K$}, muni de la structure suivante \inserted{(``$M$-épinglage'' de $K$)} : \deleted{\unsure{à|l'} \unsure{une|savoir}}

a) \deleted{\unsure{un|une} \unsure{de} \unsure{Pour}} Une application strictement croissante
\[
K\xrightarrow{\tau_K}\deleted{\unsure{I|N}}N
\qquad\text{(application ``\uline{type}'')}
\]

b) Pour tt $\xi\in K$, \inserted{posant $n=\tau_K(x)$,} une \deleted{\unsure{applic\textsuperscript{n}|applic\textsuperscript{t}|applic}} \unsure{isomorphisme|isomorphis\textsuperscript{me}}
\[
\deleted{\unsure{D|I|\varphi_\xi}}\;J_n\underset{u_\xi^K}{\simeq}K_{\leq\xi}
\qquad\text{(application d'épinglage de }K_{\leq\unsure{\mid|\xi|\eta}}\text{)}
\]
% slanted
\margin{les applications $J_n\unsure{-|\to}K$ définies par les $\xi\in K(n)=\unsure{(}\unsure{\mathrm{d.}}\tau^{-1}(n)$ sont dite \uline{admissibles}}
\[
u_\xi^K(d_n)=\xi\quad\text{(= plus grd \unsure{elément|élément|elém|élém} de }K_{\leq\xi}\text{)}
\]
avec \deleted{condition} les propriétés \deleted{de \unsure{transitivité|transit|transit.}} :
% archive sheet number: 12
